% page 497 du fac-simile (image p-496 du scan)
% bloc 165.1 x 242.0 mm ; marge gauche 36.9 mm
\pgc{15.240mm}{0.000mm}{
\l{\cel{53}489}
\l{}
\l{\sou{rekompensar}\cel{2}(\sou{trans., pro, per}). Gratifikar pro servo}
\l{\cel{3}agita. quankam ne-debata. - EFIRS.}
\l{}
\l{\sou{rekordo}. Prodajo sportala (od analoga), konstatita ofi-}
\l{\cel{3}cale, e qua ecelas omna agi antea di la sama genero.}
\l{\cel{3}- DEFR.}
\l{}
\l{\sou{rekortar}. (\sou{trans.}) Plukurtigar per tra-tranchar la extre-}
\l{\cel{3}majo (hari, kordo, e c.), per detranchar la bordo. -}
\l{}
\l{}
\l{\sou{rekrutar}. (\sou{trans.) (ulu, ad).}\cel{2}I. Obtenar la su-engajo libera (di}
\l{\cel{3}ulu) kom soldato. - II. Obtenar la su-engajo libera (diulu) a}
\l{\cel{3}trupo, a partiso, a +muvmento. - DEFIRS.}
\l{}
\l{\sou{rekta.}\cel{2}(\sou{geom.)}\cel{2}Qua, en spaco, de sua komenco til sua fino, iras}
\l{\cel{3}sen deviaco, t.e. segun la lineo maxim kurta qua esas trasebla}
\l{\cel{3}de un punto ad altra punto. - DEFIS.}
\l{}
\l{\sou{rektangulo}. (\sou{geom.)}\cel{2}Qua havas un orto. - EFIS.}
\l{}
\l{\sou{rektifikar}. (\sou{trans.}) I. (\sou{kemio}) . Igar (liquido) pura per un}
\l{\cel{3}distilo plusa. - II. (\sou{metaf}.)Igar exakta per emendo, korek-}
\l{\cel{3}tigo. - DEFIRS.}
\l{}
\l{\sou{rektolinea}. - (\sou{matem.}) Qua konsistas ye linei rekta. -}
\l{}
\l{\sou{rektoro}. La persono qua direktas universitato, distrikto univer-}
\l{\cel{3}sitatala, o Jezuito-kolegio. - DEFIRS.}
\l{}
\l{\sou{rektumoptoso}. (\sou{patol.}) Laxesko di la muskuli qui subtenas la}
\l{\cel{3}parieto di la rektumo. -}
\l{}
\l{\sou{rekuperar. (trans., de)}\cel{2}Riposedeskar (per agir por lo). - DEFIS.}
\l{}
\l{\sou{rekurenta}. I. (\sou{anat.}) (\sou{pri nervo)}. Qua kurveskas, e retrovenas}
\l{\cel{3}preske a sua departo-punto. -II. (\sou{algebro}).\cel{1}\sou{Serio rekurenta:}}
\l{\cel{3}Di qua singla ek la termini esas la sumo di ula quanto de ter-}
\l{\cel{3}mini qui preiras lu, e qui multiplikesas rispektive per ex-}
\l{\cel{3}presuri nevarianta. -}
\l{}
\l{\sou{rekursar}. (\sou{netrans.,}\cel{1}ad) I. Irar demandar (ad ulu) helpeso quan}
\l{\cel{3}onu ne recevis de altru. II. (\sou{Yurocienco)}. Apelar ad (ulu,}
\l{\cel{3}ula judicio-korto) kontre verdikto judiciala. - DEFIS.}
\l{}
\l{\sou{rekuzar}. (\sou{trans., ulu, su)}. (\sou{yurocienco)}. Refuzar aceptar kom}
\l{\cel{3}judicionto, kom juriano, kom testo, kom arbitronto. - II.}
\l{\cel{3}Refuzar agnoskar la autoritato di ulu.- III.\cel{1}\sou{rekuzar}\cel{1}su :}
\l{\cel{3}refuzar savigar sua opiniono. - FIS.}
\l{}
\l{\sou{relo}. (\sou{tekn.)}\cel{3}Singla de la du bendi ek fero o stalo, polisita,}
\l{\cel{3}sur qui la roti di treno-veturi rulas facile e mantenesas en}
\l{\cel{3}la direciono volata ; konvexan esas la parto qua kontaktas la}
\l{\cel{3}roti.(Opoze a \textquotedbl{}tramo\textquotedbl{}, ye qua ta parto esas konkava). - EFRS.}
}
